在这十五个固定通道内，TECE全流程相对PBE的三阶MAE在WS$_2$、MoS$_2$和WSe$_2$中分别为1.468、1.546和0.157，四阶分别为0.911、1.626和0.145。对应TECE Stage1 + MatterSim的三阶为6.104、4.328和1.180，四阶为2.594、2.939和0.628，全流程路线均有改善。EquiformerV3全流程的WS$_2$与MoS$_2$三阶MAE更低（0.953和0.868），但TECE在这三材料的四阶MAE更低；WSe$_2$的TECE全流程三四阶均表现较好。Prophet全流程并未普遍改善：WS$_2$四阶MAE为3.996、WSe$_2$为1.643，均高于对应Prophet+MatterSim。因此不能从Stage1表现推定其Stage2高阶精度，也不能把“大模型”合并为同一准确性类别。

固定PBE构型的WS$_2$ $\Gamma8$--M6中，QE四阶为$-1.4141$，TECE为$-1.4250$，MatterSim为$+1.9304$，Prophet为$-5.3944$，EquiformerV3为$-0.1993$ $\quartunit$。在$Q=1$时，不经拟合的偶混合原始能量分别为QE $-0.35252$、TECE $-0.35477$、MatterSim $+0.49472$、Prophet $-1.48308$、EquiformerV3 $-0.05150$ meV/超胞；$Q=0.5$同样给出TECE与QE接近的负号。由此可将此前结论收紧为\emph{MatterSim在该通道的混合非谐响应存在偏差}，而非MLFF普遍无法识别四阶符号。Prophet虽同号却过强；EquiformerV3虽在DFT基上同号却过弱，其自行弛豫模式基上接近零且为正（0.0299），提示近零通道仍对结构／模式选择敏感。

较好的四阶精度未必伴随较低的整张势能面能量MAE。例如MoS$_2$固定DFT构型的MatterSim能量MAE为42.849 meV，而TECE为46.134 meV；但TECE五通道四阶MAE为1.803，低于MatterSim的3.302。整体能量误差可能由谐波或轴向项主导，目标混合四阶仍须单独评价。表中原子力误差也作为独立指标保留。

现阶段可将TECE和EquiformerV3的Stage2作为更高精度的候选对照，MatterSim保留为快速筛选基线。该判断只基于既定五通道及当前CPU配置；本轮没有对三个新增Stage2开展全部324候选的筛选排名回放，也没有新增17$\times$17密度验证。全部数值仍是指定有限窗口的拟合系数，不等同于已证明零位移Taylor导数收敛。
